\chapter{Glossary}
\label{app:glossary}

\begin{description}
\item[Abstraction] A selective representation that omits detail for a purpose.
\item[Acceptance criterion] An observable condition used to decide whether a requirement is satisfied.
\item[Accessibility] The extent to which people with diverse abilities can perceive, operate, understand, and use a system.
\item[ACID] Atomicity, consistency, isolation, and durability properties of transactions.
\item[Actor] A role external to a system boundary that participates in an interaction.
\item[Agile] A family of principles and practices for iterative learning and delivery under change.
\item[Algorithm] A specified procedure for transforming inputs into outputs.
\item[API] An interface through which software components communicate.
\item[Authorisation] The decision about what an authenticated actor may do.
\item[Authentication] Establishing or asserting the identity of an actor.
\item[Bias] A systematic distortion or unequal pattern introduced by data, measurement, design, or use.
\item[Class] A definition from which objects with related structure or behaviour may be created.
\item[Client] A component that requests a service from another component.
\item[Complexity] A description of resources such as time or space as input size changes.
\item[Constraint] A condition that limits or validates possible states or actions.
\item[Construct] An abstract concept represented through measurement or design.
\item[Coupling] The degree to which components depend on one another.
\item[Data] Recorded distinctions or representations interpreted in a context.
\item[Database] A managed system for storing, querying, and updating structured information.
\item[Digitalisation] Organisational or social change enabled by digital technologies.
\item[Encapsulation] Protecting representation behind a controlled interface.
\item[Evaluation] A systematic judgement of an artefact, design, process, or claim against criteria.
\item[Exception] A mechanism for transferring control when normal execution cannot continue.
\item[Function] A mapping from inputs to outputs; in programming it may also have effects.
\item[Graph] A structure of vertices and edges.
\item[HCI] Human--computer interaction: study and design of interactive systems in use.
\item[HTTP] A protocol for communication between clients and servers on the web.
\item[Information system] A sociotechnical arrangement that produces and uses information.
\item[Invariant] A proposition that remains true across specified operations or states.
\item[Iteration] Repetition that may incorporate learning or changed decisions.
\item[Key] An attribute or set of attributes used to identify a tuple.
\item[Machine learning] Methods that estimate patterns or functions from data.
\item[Model] A selective representation of a phenomenon, system, or design.
\item[Normalisation] Relational decomposition guided by dependencies and anomaly reduction.
\item[Object] An entity with identity, state, and behaviour in a program.
\item[Polymorphism] The ability to use a common interface with varying implementations.
\item[Prototype] An artefact made to explore or communicate a design.
\item[Requirement] A needed capability, condition, or constraint in context.
\item[Research contribution] A defensible addition to knowledge, method, theory, design, or evidence.
\item[REST] An architectural style for networked systems based on constraints including representations and stateless interactions.
\item[Schema] A formal description of data structure and constraints.
\item[Security] Protection of assets against threats under specified assumptions.
\item[Stakeholder] A person, group, or institution affected by or able to affect a system.
\item[State] The information needed to describe a system at a time for a purpose.
\item[Technical debt] A present design choice that creates future cost or reduced options.
\item[Transaction] A coordinated unit of database work with defined consistency behaviour.
\item[Usability] Effectiveness, efficiency, and satisfaction in a specified context of use.
\item[Validation] Judging whether a model or artefact is adequate for a purpose.
\item[Verification] Checking whether an implementation conforms to a specification.
\end{description}

